
%%%===========================================
\section{Timeline}

According to the official academic calendar of the doctoral program, the student's \textbf{enrollment began on January 25, 2024}, which also marks the starting point for the time count toward graduation requirements. Thus, \textbf{the maximum allowed period for submitting the final dissertation is October 25, 2028}, with the \textbf{deadline for registering for the qualifying examination set for June 1, 2026}. Within this timeframe, the research plan is structured around the three main articles/chapters of the thesis: \textbf{Article 1}, devoted to the theoretical framework and preliminary formalization; \textbf{Article 2}, devoted to the dynamic formal model; and \textbf{Article 3}, devoted to the empirical application to Mato Grosso.

The schedule below reflects the current stage of the research. The initial formulation of the research problem, hypotheses, and conceptual architecture has already been developed for the qualifying examination. The next stages are organized to consolidate Article 1 first, then deepen the mathematical structure in Article 2, and finally implement the empirical strategy in Article 3.

\begin{table}[H]
\centering
\renewcommand{\arraystretch}{1.35}
\footnotesize
\begin{tabularx}{\textwidth}{|c|X|}
\hline
\textbf{Period} & \textbf{Milestones and Planned Activities} \\
\hline

\textbf{1st Semester 2026} 
& 
\textbf{Qualifying examination and consolidation of the research proposal.} 
Final refinement of the research problem, hypotheses, objectives, and methodological strategy. 
Organization of the thesis around three articles/chapters. 
Preliminary articulation of the concept of \emph{economic commensalism}. 
Definition of Article 1 as the immediate priority after qualification. \\
\hline

\textbf{2nd Semester 2026} 
& 
\textbf{Article 1 / Chapter 1: Theoretical Framework and Preliminary Formalization.} 
Completion of the literature review on production theory, environmental economics, common-pool resources, institutional approaches, and ecological-economic interaction regimes. 
Development of the conceptual distinction between predation, extraction, extractivism, strict preservation, and economic commensalism. 
Preliminary formal representation of the environment as an endogenous cofactor of production. 
Drafting and revision of Article 1, with the goal of producing a first complete version suitable for advisor feedback and possible submission. \\
\hline

\textbf{1st Semester 2027} 
& 
\textbf{Article 2 / Chapter 2: Dynamic Formal Model.} 
Development of the dynamic mathematical model based on the preliminary formalization introduced in Article 1. 
Specification of agents, state variables, control variables, environmental quality, payoff structure, and constraints. 
Formal treatment of intertemporal choice, uncertainty, bounded rationality, ecological feedback, and strategic interaction. 
Initial assessment of existence, stability, and equilibrium conditions. \\
\hline

\textbf{2nd Semester 2027} 
& 
\textbf{Article 2 refinement and transition to Article 3.} 
Computational simulation of incentive regimes, including positive incentives, negative incentives, mixed regimes, local compensation, and global compensation. 
Comparison between commensalist, extractivist, and predatory baselines. 
Revision of Article 2 based on formal results. 
Definition of the empirical strategy for Article 3, including data sources, spatial scale, temporal scale, treatment variables, control groups, and identification strategy. \\
\hline

\textbf{1st Semester 2028} 
& 
\textbf{Article 3 / Chapter 3: Empirical Application to Mato Grosso.} 
Collection, cleaning, and integration of economic, environmental, geospatial, and institutional data. 
Construction of indicators related to agricultural productivity, land use, deforestation, vegetation cover, environmental compliance, credit, market access, and policy exposure. 
Implementation of the empirical strategy, potentially including panel data models, spatial econometrics, difference-in-differences, propensity score matching, or synthetic control, depending on data availability and identification feasibility. \\
\hline

\textbf{2nd Semester 2028} 
& 
\textbf{Final integration, revision, and dissertation submission.} 
Consolidation of the three articles/chapters into the final dissertation. 
Cross-chapter harmonization of concepts, notation, hypotheses, and empirical interpretation. 
Final revision of theoretical claims, formal results, empirical findings, limitations, and policy implications. 
Formatting, advisor review, committee adjustments, and submission of the final dissertation before the program deadline of October 25, 2028. \\
\hline

\end{tabularx}
\end{table}

This schedule was designed to meet the formal deadlines of the doctoral program while preserving enough flexibility for theoretical refinement, mathematical development, data limitations, and empirical adjustments. The sequence is cumulative: Article 1 establishes the conceptual and preliminary formal foundations; Article 2 deepens these foundations through dynamic modeling; and Article 3 evaluates the empirical plausibility of the proposed framework in Mato Grosso. This structure also allows partial outputs to be developed as standalone academic articles while remaining integrated into the broader thesis.

\newpage

% fim %
